\chapter{Diagnostic Decision Trees --- What Happens Next}
\label{app:algorithms}

\clearpage
\section{Introduction}

When your doctor orders a test and the result is abnormal, there is usually a step-by-step process for figuring out what is going on. These decision trees --- called ``diagnostic algorithms'' --- are simplified educational schematics, not standalone clinical protocols. They must be adapted to the patient's age, pregnancy status, comorbidities, stability, and local guidelines.

This appendix walks you through the most common situations: you have an abnormal result, and now what?

\clearpage
\section{Anemia --- You Are Tired, Your Hemoglobin Is Low}

Anemia means your blood is not carrying enough oxygen. It is the most common blood disorder worldwide. The workup follows a clear path.

\subsection{Step 1: Confirm the Anemia}

Your CBC shows a low hemoglobin or hematocrit. The first question: \emph{is this real?}

\begin{itemize}
    \item If you were dehydrated when the blood was drawn, your hemoglobin may look falsely low. If you were overhydrated, it may look falsely normal.
    \item Repeat the CBC to confirm.
\end{itemize}

\subsection{Step 2: Reticulocyte Count --- Is Your Bone Marrow Responding?}

The reticulocyte count tells your doctor whether your bone marrow is trying to fix the problem.

\begin{itemize}
    \item \textbf{High reticulocyte count:} Your bone marrow is working hard. The problem is either blood loss (you are bleeding somewhere) or hemolysis (your red blood cells are being destroyed). Go to Step 3a.
    \item \textbf{Low or normal reticulocyte count:} Your bone marrow is not responding. The problem is that it does not have the raw materials it needs, or it is suppressed. Go to Step 3b.
\end{itemize}

\subsection{Step 3a: If the Bone Marrow Is Responding (High Reticulocyte Count)}

If you are making lots of new red blood cells, the problem is that something is destroying or losing them:

\begin{itemize}
    \item \textbf{Check for bleeding:} Stool occult blood test (are you bleeding from your gut?), urine analysis (blood in urine?), menstrual history (heavy periods?), and in older adults, colonoscopy to rule out colon cancer.
    \item \textbf{Check for hemolysis:} LDH (elevated when red blood cells break), haptoglobin (low when hemolysis is happening), indirect bilirubin (elevated from red blood cell breakdown), reticulocyte count (high).
    \item \textbf{Peripheral blood smear:} A pathologist looks at your blood under a microscope. Abnormal red blood cell shapes can point to specific causes (sickle cells, schistocytes from DIC, spherocytes from autoimmune hemolysis).
    \item \textbf{Direct Coombs test:} Tests whether your immune system is attacking your own red blood cells (autoimmune hemolytic anemia).
\end{itemize}

\subsection{Step 3b: If the Bone Marrow Is Not Responding (Low Reticulocyte Count)}

If your bone marrow is not making enough new red blood cells, you need to figure out why. The MCV (size of your red blood cells) is the key starting point.

\subsubsection{Small Red Blood Cells (MCV $<$80 fL)}

\begin{itemize}
    \item \textbf{Iron studies:} Serum iron (low in iron deficiency), TIBC (high in iron deficiency), ferritin (the most reliable test for iron stores --- low in iron deficiency).
    \item \textbf{If iron is low and ferritin is low:} Iron deficiency anemia. Find the cause --- most commonly slow chronic blood loss (heavy periods, GI bleeding from ulcers, colon cancer).
    \item \textbf{If iron is low but ferritin is normal or high:} Anemia of chronic disease. The iron is locked away in storage because of inflammation. Treat the underlying disease.
    \item \textbf{Hemoglobin electrophoresis:} If iron studies are normal but the MCV is still low, check for thalassemia trait. This is a genetic condition common in Mediterranean, African, and Asian populations.
\end{itemize}

\subsubsection{Normal-Sized Red Blood Cells (MCV 80--100 fL)}

\begin{itemize}
    \item \textbf{Check kidney function:} Creatinine and eGFR. Kidney disease is a common cause of normocytic anemia because your kidneys produce erythropoietin (EPO), the hormone that stimulates red blood cell production.
    \item \textbf{Check for chronic disease:} CRP, ESR, and clinical history. Chronic infections, autoimmune diseases, and cancer can all suppress bone marrow.
    \item \textbf{Check for hemolysis:} LDH, haptoglobin, bilirubin.
    \item \textbf{Consider bone marrow biopsy:} If no cause is found, a bone marrow biopsy may be needed to evaluate for aplastic anemia, myelodysplastic syndrome, or marrow infiltration.
\end{itemize}

\subsubsection{Large Red Blood Cells (MCV $>$100 fL)}

\begin{itemize}
    \item \textbf{Vitamin B12 level:} Low in pernicious anemia, malabsorption, strict veganism, or ileal disease.
    \item \textbf{Folate level:} Low from poor diet, alcohol abuse, pregnancy, or certain medications (methotrexate, phenytoin).
    \item \textbf{TSH:} Hypothyroidism can cause macrocytic anemia.
    \item \textbf{Liver function tests:} Liver disease can cause macrocytosis.
    \item \textbf{Consider myelodysplastic syndrome:} If B12, folate, and thyroid function are all normal, a bone marrow biopsy may be warranted, especially in older adults.
\end{itemize}

\subsection{Understanding Your Results}

Anemia is a symptom, not a diagnosis. Your doctor must find the underlying cause. Do not accept ``you are a little anemic'' without a workup. Ask:
\begin{itemize}
    \item ``What type of anemia is this?''
    \item ``What is causing it?''
    \item ``Do I need iron supplements, or is something else going on?''
    \item ``Should I see a gastroenterologist to look for bleeding?''
    \item ``Do I need a hematologist?''
\end{itemize}

\clearpage
\section{Chest Pain --- Is It a Heart Attack?}

Chest pain is one of the most common reasons people go to the emergency room. The workup is designed to rule in or rule out acute coronary syndrome (ACS) --- the medical term for a heart attack or its precursors.

\subsection{Step 1: ECG Within 10 Minutes}

The moment you arrive with chest pain, an ECG (electrocardiogram) is performed. This takes 5 minutes and can immediately show if you are having a STEMI --- a type of heart attack where a coronary artery is completely blocked. STEMI requires emergency cardiac catheterization.

\subsection{Step 2: Troponin at 0 Hours}

Troponin is the gold standard blood test for detecting heart muscle damage. It is drawn immediately and repeated at 1--3 hours depending on the protocol.

\begin{itemize}
    \item \textbf{Troponin negative, ECG normal:} ACS is unlikely but not completely ruled out. Continue observation and repeat troponin.
    \item \textbf{Troponin positive:} Heart muscle damage is occurring. The cardiologist is called. You are going to the catheterization lab.
\end{itemize}

\subsection{Step 3: Risk Stratification}

If the initial evaluation is not definitive, your doctor uses scoring systems to estimate your risk:

\begin{itemize}
    \item \textbf{HEART score:} History, ECG, Age, Risk factors, Troponin. A high score means high risk and you are admitted. A low score means you may be safe to go home with outpatient follow-up.
    \item \textbf{TIMI score:} Used for unstable angina and NSTEMI to guide treatment decisions.
    \item \textbf{GRACE score:} Predicts in-hospital and 6-month mortality in ACS.
\end{itemize}

\subsection{Step 4: If ACS Is Confirmed}

\begin{itemize}
    \item Antiplatelet therapy (aspirin + clopidogrel/ticagrelor)
    \item Anticoagulation (heparin)
    \item Cardiology consultation
    \item Coronary angiography and possible stenting
    \item If complete blockage: emergency PCI (percutaneous coronary intervention) within 90 minutes
\end{itemize}

\subsection{Step 5: If ACS Is Ruled Out}

Chest pain has many other causes. Your doctor will evaluate for:
\begin{itemize}
    \item \textbf{Pulmonary embolism:} Blood clot in the lungs. D-dimer, CT pulmonary angiography.
    \item \textbf{Aortic dissection:} Tearing of the aorta. CT angiography. D-dimer may be elevated.
    \item \textbf{Pneumonia:} Chest X-ray, WBC, CRP.
    \item \textbf{GERD (acid reflux):} The most common non-emergency cause of chest pain.
    \item \textbf{Musculoskeletal:} Costochondritis (inflammation of the cartilage connecting ribs to the breastbone).
    \item \textbf{Panic disorder:} A diagnosis of exclusion --- only after serious causes are ruled out.
\end{itemize}

\subsection{Understanding Your Results}

If you have chest pain, \textbf{do not drive yourself to the hospital.} Call 911. Paramedics can start an ECG in the ambulance, and if you are having a STEMI, every minute of delay increases the amount of heart muscle that dies. If you are having a heart attack, you want to be in an ambulance with trained personnel, not stuck in traffic.

\clearpage
\section{Diabetes --- How Do You Know If You Have It?}

Diabetes screening is one of the most important preventive health measures. Here is how the diagnosis works.

\subsection{Understanding Your Results}

\begin{longtable}{@{}lll@{}}
\toprule
\textbf{Test} & \textbf{Diabetes} & \textbf{Prediabetes} \\
\midrule
Fasting glucose & $\geq$126 mg/dL & 100--125 mg/dL \\
2-hour oral glucose tolerance test (OGTT) & $\geq$200 mg/dL & 140--199 mg/dL \\
HbA1c & $\geq$6.5\% & 5.7--6.4\% \\
Random glucose & $\geq$200 mg/dL with symptoms & --- \\
\bottomrule
\end{longtable}

\subsection{Understanding Your Results}

\begin{enumerate}
    \item \textbf{Fasting glucose 70--99:} Normal. Rescreen in 3 years if you have risk factors (family history, obesity, sedentary lifestyle, certain ethnicities).
    \item \textbf{Fasting glucose 100--125 (prediabetes):} Lifestyle changes (weight loss, exercise) are the first-line treatment. Repeat HbA1c in 6--12 months.
    \item \textbf{Fasting glucose $\geq$126:} Repeat on a separate day to confirm. If confirmed, you have diabetes.
    \item \textbf{HbA1c 5.7--6.4\%:} Prediabetes. Same as above.
    \item \textbf{HbA1c $\geq$6.5\%:} Diabetes. Confirm with a repeat test if you have no symptoms. If you have symptoms (excessive thirst, frequent urination, weight loss), one test is sufficient.
    \item \textbf{Random glucose $\geq$200 with symptoms:} Diabetes. No repeat test needed.
\end{enumerate}

\subsection{Understanding Your Results}

If you are diagnosed with prediabetes, \textbf{this is a wake-up call, not a death sentence.} The Diabetes Prevention Program study showed that lifestyle intervention (7\% weight loss + 150 minutes of exercise per week) reduced the progression to diabetes by 58\% --- more effective than metformin.

If you are diagnosed with diabetes, ask your doctor:
\begin{itemize}
    \item ``What is my HbA1c and what is my target?''
    \item ``Do I need medication, or can I manage this with lifestyle changes?''
    \item ``Should I see an endocrinologist?''
    \item ``How often should I have my eyes examined?'' (diabetic retinopathy screening)
    \item ``How often should I have my feet checked?'' (diabetic neuropathy)
    \item ``What is my kidney function?'' (diabetic nephropathy screening)
\end{itemize}

\clearpage
\section{Thyroid Dysfunction --- Is Your Thyroid Working?}

\subsection{Step 1: TSH Is the Screening Test}

TSH is the single best test for thyroid function. If your doctor suspects thyroid disease, TSH is drawn first.

\subsection{Step 2: Interpret the Pattern}

\begin{longtable}{@{}p{3.5cm}p{3.0cm}p{2.0cm}p{5.5cm}@{}}
\toprule
\textbf{TSH} & \textbf{Free T4} & \textbf{Diagnosis} & \textbf{What Happens Next} \\
\midrule
High & Low & Overt hypothyroidism & Start levothyroxine replacement \\
High & Normal & Subclinical hypothyroidism & Monitor or treat based on symptoms \\
Normal & Normal & Euthyroid (normal) & No thyroid treatment needed \\
Low & High & Overt hyperthyroidism & Endocrinology referral, treatment options \\
Low & Normal & Subclinical hyperthyroidism & Monitor, check for atrial fibrillation, consider treatment if elderly \\
\bottomrule
\end{longtable}

\subsection{Step 3: If the Pattern Is Unclear}

\begin{itemize}
    \item \textbf{Thyroid antibodies (TPO, TgAb):} If TSH is elevated, these antibodies suggest Hashimoto's thyroiditis (autoimmune hypothyroidism) --- the most common cause of hypothyroidism in iodine-sufficient regions.
    \item \textbf{TRAb (TSH receptor antibodies):} If TSH is low and T4 is high, these antibodies suggest Graves' disease (autoimmune hyperthyroidism).
    \item \textbf{Thyroid ultrasound:} If a nodule is felt on examination or if hyperthyroidism is confirmed.
    \item \textbf{Radioactive iodine uptake:} Distinguishes between different causes of hyperthyroidism.
\end{itemize}

\subsection{Understanding Your Results}

Hypothyroidism is extremely common, especially in women over 50. The treatment (levothyroxine) is simple, inexpensive, and well-tolerated. It must be taken on an empty stomach, 30--60 minutes before breakfast, and should not be taken with calcium, iron, or antacids. Your TSH is checked 6--8 weeks after starting or adjusting the dose.

Hyperthyroidism is less common but potentially more dangerous. Untreated hyperthyroidism can cause heart rhythm disorders (atrial fibrillation), bone loss, and in rare cases, thyroid storm (a life-threatening escalation).

\clearpage
\section{Liver Function Abnormalities --- What Pattern Do You See?}

\subsection{Step 1: Identify the Pattern}

\begin{longtable}{@{}p{4.5cm}p{10.5cm}@{}}
\toprule
\textbf{Pattern} & \textbf{What It Means} \\
\midrule
AST and ALT very high ($>$5$\times$ upper limit) & Liver cells are being damaged \\
ALP and GGT very high ($>$3$\times$ upper limit) & Bile flow is blocked \\
AST and ALT mildly elevated (1--3$\times$ upper limit) & Many possible causes \\
Isolated GGT elevation & Alcohol or drug enzyme induction \\
Isolated bilirubin elevation & Gilbert syndrome (benign) or hemolysis \\
\bottomrule
\end{longtable}

\subsection{Step 2: Follow the Pattern}

\subsubsection{If AST and ALT Are Very High (Hepatocellular Pattern)}

\begin{itemize}
    \item \textbf{Viral hepatitis:} Hepatitis A, B, C serology. Hepatitis A is from contaminated food/water. Hepatitis B and C are bloodborne.
    \item \textbf{Drug toxicity:} Acetaminophen overdose (check level), herbal supplements (kava, comfrey), prescription medications (statins, isoniazid, valproic acid).
    \item \textbf{Alcoholic hepatitis:} History of heavy drinking, AST:ALT ratio typically $>$2:1.
    \item \textbf{Ischemic hepatitis:} Low blood flow to the liver (heart failure, shock). AST and ALT can exceed 1000 U/L.
    \item \textbf{Autoimmune hepatitis:} ANA, anti-smooth muscle antibodies.
    \item \textbf{Wilson disease:} In young patients. Ceruloplasmin, copper studies.
\end{itemize}

\subsubsection{If ALP and GGT Are Very High (Cholestatic Pattern)}

\begin{itemize}
    \item \textbf{Gallstones:} Abdominal ultrasound is the first test.
    \item \textbf{Bile duct obstruction:} MRCP (magnetic resonance cholangiopancreatography) to visualize the bile ducts.
    \item \textbf{Primary biliary cholangitis:} Anti-mitochondrial antibodies (AMA). More common in middle-aged women.
    \item \textbf{Primary sclerosing cholangitis:} Associated with inflammatory bowel disease. MRCP shows ``beading'' of bile ducts.
    \item \textbf{Drug-induced cholestasis:} Check medication list.
    \item \textbf{Infiltrative disease:} Sarcoidosis, lymphoma, amyloidosis.
\end{itemize}

\subsection{Understanding Your Results}

Mildly elevated liver enzymes (1--3$\times$ normal) are extremely common and are often caused by fatty liver disease (non-alcoholic fatty liver disease, or NAFLD), obesity, alcohol, and medications. Your doctor should investigate, but do not panic. The most important thing you can do for your liver is maintain a healthy weight, limit alcohol, and avoid unnecessary medications and supplements.

\clearpage
\section{Thyroid Nodules --- You Have a Lump in Your Thyroid}

\subsection{Step 1: Ultrasound}

When a thyroid nodule is found (on physical exam or incidentally on imaging), thyroid ultrasound is performed to characterize it.

\subsection{Step 2: TI-RADS Scoring}

Thyroid nodules are scored using the ACR TI-RADS system based on their ultrasound appearance:
\begin{itemize}
    \item Composition (solid vs. cystic)
    \item Echogenicity (how bright it is)
    \item Shape (wider than tall or taller than wide)
    \item Margins (smooth, irregular, or lobulated)
    \item Echogenic foci (calcifications)
\end{itemize}

Higher TI-RADS scores require fine needle aspiration (FNA) biopsy.

\subsection{Step 3: FNA Biopsy}

If the nodule is large enough and suspicious enough, a thin needle is inserted into the nodule under ultrasound guidance to collect cells for examination. Results are reported using the Bethesda system:
\begin{itemize}
    \item \textbf{Benign:} No cancer. Follow with ultrasound.
    \item \textbf{Malignant:} Cancer confirmed. Surgery.
    \item \textbf{Indeterminate:} Cannot determine. May need molecular testing or diagnostic surgery.
\end{itemize}

\subsection{Understanding Your Results}

Most thyroid nodules are benign (85--95\%). Thyroid cancer, when it occurs, is usually slow-growing and highly treatable. Papillary thyroid carcinoma, the most common type, has a 20-year survival rate exceeding 95\%. Do not let the word ``cancer'' send you into panic --- thyroid cancer is one of the most curable cancers.

\clearpage
\section{Sepsis --- Could This Be a Life-Threatening Infection?}

Sepsis is your body's overwhelming, dysregulated response to infection. It is the leading cause of death in hospitals.

\subsection{Step 1: Immediate Clinical Assessment}

Assess for suspected infection, abnormal vital signs, altered mental status, hypoxemia, impaired perfusion, and acute organ dysfunction. Do not use qSOFA as the sole sepsis screening tool: it is not sufficiently sensitive to identify all patients at risk. Hospitals may use an internally validated early-warning or sepsis-screening system, but screening must be combined with clinical judgment and prompt reassessment.

qSOFA can be used as a risk-stratification prompt in patients with suspected infection:
\begin{itemize}
    \item Respiratory rate $\geq$22 breaths per minute
    \item Altered mental status (Glasgow Coma Scale $<$15)
    \item Systolic blood pressure $\leq$100 mmHg
\end{itemize}

Two or more criteria identify a patient at higher risk of poor outcome and should prompt urgent assessment; qSOFA does not diagnose sepsis and a score below two does not exclude it.

\subsection{Step 2: Confirm and Evaluate}

\begin{itemize}
    \item \textbf{Lactate level:} An elevated lactate can indicate impaired perfusion or other metabolic stress, but it is nonspecific. Use serial values and the clinical context; modern sepsis definitions do not classify sepsis solely by a lactate threshold.
    \item \textbf{Blood cultures (at least 2 sets):} Drawn from two separate sites before antibiotics are started. These identify the specific bacteria causing the infection.
    \item \textbf{Complete blood count:} WBC may be very high, very low, or normal in sepsis. Bandemia (immature neutrophils) is a red flag.
    \item \textbf{Procalcitonin:} A biomarker that rises in bacterial infections and helps distinguish bacterial from viral.
    \item \textbf{Organ function tests:} Creatinine (kidney), bilirubin (liver), platelets (coagulation), troponin (heart).
\end{itemize}

\subsection{Step 3: Treatment (The Hour Matters)}

For septic shock or a high likelihood of sepsis, begin resuscitation and antimicrobial treatment urgently according to the current local protocol. The following are common guideline-based actions, not automatic orders for every patient:
\begin{itemize}
    \item Broad-spectrum antibiotics
    \item IV crystalloid for sepsis-induced hypoperfusion or septic shock, with repeated assessment for fluid responsiveness and overload; the volume and rate must be individualized
    \item Vasopressors if blood pressure does not respond to fluids
    \item Source control (drain abscess, remove infected device)
\end{itemize}

\subsection{Understanding Your Results}

Sepsis can be a medical emergency. Confusion, rapid breathing, low blood pressure, severe breathlessness, reduced urine output, mottled or blue skin, or rapidly worsening illness require urgent medical assessment. Treatment timing depends on the likelihood of sepsis and whether shock is present; avoid presenting a universal hourly mortality estimate as a decision rule.

\clearpage
\section{Hypoglycemia --- Why Is Your Blood Sugar Dangerously Low?}

\subsection{Step 1: Confirm True Hypoglycemia}

True hypoglycemia requires Whipple's triad:
\begin{enumerate}
    \item Symptoms consistent with low blood sugar (shakiness, sweating, confusion, hunger)
    \item Measured glucose $<$70 mg/dL at the time of symptoms
    \item Relief of symptoms when glucose is raised
\end{enumerate}

If all three criteria are met, you have true hypoglycemia and it needs to be investigated.

\subsection{Step 2: Check Insulin and C-Peptide During the Episode}

The key to diagnosing the cause is checking insulin, C-peptide, and proinsulin during the hypoglycemic episode.

\begin{longtable}{@{}p{2.0cm}p{2.5cm}p{9.0cm}@{}}
\toprule
\textbf{Insulin} & \textbf{C-Peptide} & \textbf{Diagnosis} \\
\midrule
High & High & Insulinoma (insulin-producing tumor) or sulfonylurea use \\
High & Low & Exogenous insulin injection (someone is injecting insulin) \\
High & High & Sulfonylurea ingestion (check drug levels) \\
Low & Low & Non-islet cell tumor, adrenal insufficiency, severe illness \\
\bottomrule
\end{longtable}

\subsection{Understanding Your Results}

Recurrent hypoglycemia is not normal. If you are not diabetic and you keep having low blood sugar episodes, you need a thorough workup. The most concerning cause is insulinoma --- a small tumor of the pancreas that produces insulin autonomously. It is usually benign and surgically curable.

If you are diabetic and having recurrent hypoglycemia, your medication dose likely needs adjustment. Tell your doctor --- do not just eat more sugar to compensate. Recurrent hypoglycemia, especially in older adults, is associated with increased risk of falls, fractures, and cognitive decline.

\clearpage
\section{How to Use These Decision Trees}

These decision trees are simplified educational versions of common diagnostic reasoning. They are not a substitute for professional medical judgment, local laboratory policy, or current specialty guidance. They help you understand:

\begin{itemize}
    \item \textbf{Why your doctor is ordering more tests.} A single test is rarely the end of the story.
    \item \textbf{Why some results are repeated.} Confirmation is a standard part of good medicine.
    \item \textbf{Why timing matters.} Some tests must be drawn at specific times or during specific symptoms to be meaningful.
    \item \textbf{When to push for more answers.} If your doctor says ``your tests are fine'' but you still have symptoms, use this knowledge to ask informed follow-up questions.
    \item \textbf{When to seek emergency care.} Some patterns require immediate intervention.
\end{itemize}
